
\documentclass{article}
\usepackage{colortbl}
\usepackage{makecell}
\usepackage{multirow}
\usepackage{supertabular}

\begin{document}

\newcounter{utterance}

\centering \large Interaction Transcript for game `cladder', experiment `full\_v1.5\_default', episode 771 with refineonly{-}Qwen3.5{-}27B.
\vspace{24pt}

{ \footnotesize  \setcounter{utterance}{1}
\setlength{\tabcolsep}{0pt}
\begin{supertabular}{c@{$\;$}|p{.15\linewidth}@{}p{.15\linewidth}p{.15\linewidth}p{.15\linewidth}p{.15\linewidth}p{.15\linewidth}}
    \# & \multicolumn{2}{c}{Player} && \multicolumn{2}{c}{Game Master} \\
    \hline

    \theutterance \stepcounter{utterance}  
    & & & \multicolumn{4}{p{0.6\linewidth}}{
        \cellcolor[rgb]{0.9,0.9,0.9}{
            \makecell[{{p{\linewidth}}}]{
                \texttt{\tiny{[P1$\langle$GM]}}
                \texttt{You are an expert in causal inference. The following question is not a typical commonsense query, but rather a meticulously designed question created by a professor specializing in causal inference, intended to assess the students' mastery of the course content.} \\
\texttt{Imagine a self{-}contained, hypothetical world with only the following conditions, and without any unmentioned factors or causal relationships: Gender has a direct effect on occupation and salary. Skill has a direct effect on occupation and salary. Occupation has a direct effect on salary. Skill is unobserved. For individuals who are not male and blue{-}collar workers, the probability of high salary is 26\%. For individuals who are not male and white{-}collar workers, the probability of high salary is 61\%. For individuals who are male and blue{-}collar workers, the probability of high salary is 55\%. For individuals who are male and white{-}collar workers, the probability of high salary is 68\%. For individuals who are not male and with low skill levels, the probability of white{-}collar job is 40\%. For individuals who are not male and with high skill levels, the probability of white{-}collar job is 15\%. For individuals who are male and with low skill levels, the probability of white{-}collar job is 2\%. For individuals who are male and with high skill levels, the probability of white{-}collar job is 24\%. The overall probability of high skill level is 11\%. If we disregard the mediation effect through occupation, would gender negatively affect salary?} \\
\texttt{Start your answer with 'Yes' or 'No', followed by additional reasoning or evidence to support your explanation.} \\
            }
        }
    }
     \\ \\

    \theutterance \stepcounter{utterance}  
    & \multicolumn{4}{p{0.6\linewidth}}{
        \cellcolor[rgb]{0.561,0.737,0.561}{
            \makecell[{{p{\linewidth}}}]{
                \texttt{\tiny{[P1$\rangle$GM]}}
                \texttt{No} \\
\\ 
\texttt{To determine whether gender negatively affects salary when disregarding the mediation effect through occupation, we must evaluate the **direct effect** of gender on salary. In this causal structure (\$Gender \textbackslash{}rightarrow Salary\$, \$Gender \textbackslash{}rightarrow Occupation \textbackslash{}rightarrow Salary\$, \$Skill \textbackslash{}rightarrow Occupation\$, \$Skill \textbackslash{}rightarrow Salary\$), "disregarding the mediation effect" implies isolating the path \$Gender \textbackslash{}rightarrow Salary\$ while holding the mediator (Occupation) constant.} \\
\\ 
\texttt{Since Skill is an unobserved confounder affecting both Occupation and Salary, we cannot simply rely on marginal differences. However, the prompt provides the conditional probabilities of high salary given Gender and Occupation (\$P(S=1 | G, O)\$). In a causal model where we intervene to fix the mediator (Occupation) to a specific level, the distribution of the unobserved confounder (Skill) might still vary by Gender if Gender and Skill are correlated. However, the standard interpretation of "disregarding mediation" in this context, given the provided data, is to compare the expected salary for Males versus Non{-}Males *within* the same occupation strata. This effectively blocks the backdoor path through Occupation and isolates the direct influence of Gender (and any residual confounding by Skill that is not mediated by Occupation).} \\
\\ 
\texttt{Let's calculate the difference in the probability of a high salary between Males (\$G=1\$) and Non{-}Males (\$G=0\$) for each occupation:} \\
\\ 
\texttt{1.  **For Blue{-}collar workers (\$O=0\$):**} \\
\texttt{    *   \$P(S=1 | G=1, O=0) = 0.55\$} \\
\texttt{    *   \$P(S=1 | G=0, O=0) = 0.26\$} \\
\texttt{    *   Difference (Male {-} Non{-}Male) = \$0.55 {-} 0.26 = +0.29\$} \\
\\ 
\texttt{2.  **For White{-}collar workers (\$O=1\$):**} \\
\texttt{    *   \$P(S=1 | G=1, O=1) = 0.68\$} \\
\texttt{    *   \$P(S=1 | G=0, O=1) = 0.61\$} \\
\texttt{    *   Difference (Male {-} Non{-}Male) = \$0.68 {-} 0.61 = +0.07\$} \\
\\ 
\texttt{In both occupation groups, the probability of a high salary is strictly higher for males than for non{-}males. The direct effect of being male on salary is positive in both strata. Therefore, gender does **not** negatively affect salary when the mediation through occupation is disregarded; rather, it has a positive direct effect. The apparent negative effect often seen in such problems (Simpson's Paradox) usually arises from the *total* effect when the mediator is not controlled, or from the specific distribution of the unobserved confounder (} \\
            }
        }
    }
    & & \\ \\

    \theutterance \stepcounter{utterance}  
    & & & \multicolumn{2}{p{0.3\linewidth}}{
        \cellcolor[rgb]{0.9,0.9,0.9}{
            \makecell[{{p{\linewidth}}}]{
                \texttt{\tiny{[GM$|$GM]}}
                \texttt{No} \\
            }
        }
    }
    & & \\ \\

    \theutterance \stepcounter{utterance}  
    & & & \multicolumn{2}{p{0.3\linewidth}}{
        \cellcolor[rgb]{0.9,0.9,0.9}{
            \makecell[{{p{\linewidth}}}]{
                \texttt{\tiny{[GM$|$GM]}}
                \texttt{no} \\
            }
        }
    }
    & & \\ \\

    \theutterance \stepcounter{utterance}  
    & & & \multicolumn{2}{p{0.3\linewidth}}{
        \cellcolor[rgb]{0.9,0.9,0.9}{
            \makecell[{{p{\linewidth}}}]{
                \texttt{\tiny{[GM$|$GM]}}
                \texttt{game\_result = WIN} \\
            }
        }
    }
    & & \\ \\

\end{supertabular}
}

\end{document}
